\section{Appendix Overview}
\label{app:overview}
The appendix is an audit trail for the main paper rather than a second results narrative. It follows the same order as the experiment: data selection, provider surfaces, response generation, oracle construction, and supporting diagnostics.

\begin{description}[leftmargin=0pt,itemsep=0.25em,topsep=0.25em]
\item[Data and providers.] Appendix~\ref{app:dataset} documents how the \textsc{SealQA-Hard} sample was selected; Appendix~\ref{app:providers} fixes the provider requests and shared page-fetch backend.
\item[Generation and traces.] Appendix~\ref{app:generation} records the agent prompt contract, trace schema, and qualitative audit views.
\item[Metrics and results.] Appendix~\ref{app:metrics} defines the Kimi per-URL oracle, validation checks, and derived metrics; Appendix~\ref{app:additional} collects uncertainty intervals, decision-cell diagnostics, token/fetch summaries, and artifact provenance.
\end{description}

\section{Dataset Construction}
\label{app:dataset}
The experiment uses a deterministic 100-query sample from \textsc{SealQA-Hard}. We assign stable query IDs by hashing question text, form strata over \texttt{freshness} $\times$ \texttt{search\_results} $\times$ \texttt{topic}, allocate slots by largest-remainder proportional rounding, and sample within cells using seed 20260509.

The table below records the concrete sample inputs and audit breadcrumbs so the sampling procedure can be checked without relying on prose alone.

\begin{center}
\footnotesize
\begin{tabularx}{\linewidth}{@{}L{0.34\linewidth}Z@{}}
\toprule
\textbf{Dataset field} & \textbf{Value} \\
\midrule
Dataset & \texttt{vtllms/sealqa} \\
Subset & \texttt{seal\_hard} \\
Source rows & 254 \\
Selected rows & \NQueries{} \\
Seed & 20260509 \\
Primary strata & \texttt{freshness}, \texttt{search\_results}, \texttt{topic} \\
Source file & \path{data/raw/seal-hard.jsonl} \\
Selected sample & \path{data/queries/phase1_100.json} \\
Sampling audit & \begin{tabular}[t]{@{}l@{}}\path{data/100-dataset-selection-}\\\path{rationale.md}\end{tabular} \\
Offline fields & Gold answers, gold URLs, metadata, and source indices are retained for grading and audit only. \\
\bottomrule
\end{tabularx}
\end{center}

The selected sample preserves the source distribution closely. Marginal counts are: 25 fast-changing, 43 slow-changing, and 32 never-changing questions; 57 conflicting and 43 unhelpful search-result labels; topics are Science \& Technology 27, Sports 22, Entertainment 22, Others 12, Politics 9, and History \& Geography 8. Effective years are 61 before 2024, 16 in 2024, and 23 in 2025. Multi-label question-type tags include 72 advanced reasoning, 61 entity/event disambiguation, 13 temporal tracking, 5 cross-lingual reasoning, and 2 false-premise questions.

\section{Provider Settings}
\label{app:providers}
The provider configuration is fixed by \path{config/experiment.yaml} and the provider-specific YAML files under \path{config/providers/}. All providers return web-search results rather than provider-side page extracts, and each adapter normalizes output to the same model-visible schema before rendering.

\paragraph{Brave.}
The adapter issues a GET request to \path{/res/v1/web/search} with count 10, offset 0, US/en locale, \texttt{safesearch=moderate}, \texttt{result\_filter=web}, and text decorations disabled. The rendered snippet surface uses \texttt{description} plus provider-native \texttt{extra\_snippets}.

\paragraph{Tavily.}
The adapter issues a POST request to \path{/search} with \texttt{search\_depth=basic}, \texttt{topic=general}, max 10, answer/images/favicon disabled, and \texttt{include\_raw\_content=false}. The rendered snippet surface uses \texttt{content}.

\paragraph{Firecrawl.}
The adapter issues a POST request to \path{/v2/search} with limit 10, web source only, US region, 60s provider timeout, invalid URLs retained, and provider-side markdown scraping disabled. The rendered snippet surface uses \texttt{description} or metadata description.

The normalized result fields are \texttt{rank}, \texttt{title}, \texttt{url}, \texttt{domain}, \texttt{snippet}, and \texttt{provider\_metadata}. URL normalization removes fragments and common tracking parameters. Brave's additional snippet blocks are rendered as \texttt{<extra\_snippet>} elements because they were visible to the agent as part of the provider surface; we do not reclassify them as fetched-page evidence.

\subsection{Shared Page Fetch Backend}
\label{app:fetch}
All page text in the main experiment comes from the same \fetchPage{} backend, independent of the search provider. The backend is Jina Reader markdown through \path{r.jina.ai}; it is invoked only after the agent selects a prior \texttt{document\_id}. Fetch artifacts are cached as gzip JSON records keyed by normalized URL and backend. Each record stores fetch status, final URL, content type, truncation flag, extractor name, character/token counts, a text hash, latency, and the extracted body. The configured guardrails are a 15s timeout, a 2MB maximum read, and concurrency 4. In the analyzed run, successful fetches are recorded as \texttt{jina\_reader\_markdown}; failures are retained as failed fetches rather than silently dropped.

\section{Response Generation}
\label{app:generation}
\label{app:agent-contract}
The answer model, prompt, tools, and loop budget are fixed across providers. Trace rows preserve the exact rendered request snapshots in \path{iterations[].llm_request.messages}; the prompt templates live under \path{src/searchapi_eval/agent/prompts/}.

The next table lists the fixed generation contract. Its purpose is to show which parts of the agent were held constant while the search provider changed.

\begin{center}
\footnotesize
\begin{tabularx}{\linewidth}{@{}L{0.32\linewidth}Z@{}}
\toprule
\textbf{Component} & \textbf{Fixed value} \\
\midrule
Answer model & GPT-5.4 Azure deployment, temperature 0 \\
Prompt templates & \path{system_fetch_tool.liquid}, \path{user_query.liquid}, \path{search_documents.liquid} \\
Loop budget & 10 iterations, up to 10 results per search call \\
Tools & \searchWeb{} over one configured provider; \fetchPage{} using a model-visible \texttt{document\_id} \\
Turn structure & The model may issue multiple search or fetch calls in one turn; each call is traced separately. \\
Search observation & \texttt{document\_id}, rank, title, URL, domain, snippet surface, and provider metadata \\
Fetch observation & Source provenance, fetch metadata, and extracted markdown/text from the shared backend \\
Final answer rule & Response must begin with \texttt{FINAL ANSWER:} and contain only the short factual answer. \\
Hidden fields & Gold answers and gold URLs are retained in traces for offline grading, never sent to the answer model. \\
\bottomrule
\end{tabularx}
\end{center}

In fetch-tool mode, \searchWeb{} observations are snippet-only. Page text enters the model context only after the model chooses a prior \texttt{document\_id}. This makes fetch behavior observable and keeps provider comparisons from reflecting automatic page inclusion.

\subsection{Trace Schema and Prompt Artifacts}
\label{app:trace-repro}
Each provider-query run is one JSONL trace row using the schema in \path{data/trace-schema-v1.md}. The trace records LLM request snapshots, model responses, tool calls, normalized search responses, raw provider payloads, fetch records, token counts, latency, final answers, and offline gold fields. The model-visible search surface is rendered as \texttt{<search\_documents>} with one \texttt{<document>} block per result. Fetched pages appear as \texttt{<fetched\_page>} blocks with URL provenance, fetch metadata, and extracted body content.

The artifact map below documents the internal trace schema and where information enters the audit pipeline; it does not imply redistribution of the recorded content.

\begin{center}
\footnotesize
\begin{tabularx}{\linewidth}{@{}L{0.36\linewidth}Z@{}}
\toprule
\textbf{Artifact} & \textbf{Where it appears} \\
\midrule
Rendered answer prompt & \path{iterations[].llm_request.messages} \\
Tool schemas & \path{iterations[].llm_request.tools} \\
Normalized search response & \path{retrievals[].search_response.results[]} \\
Raw provider response & \path{retrievals[].search_response.raw_response} in recorded traces \\
Agent fetch decision & \path{fetches[].requested_document_id} and \path{fetches[].reason} \\
Fetched page result & \path{fetches[].page_fetch} and the rendered \texttt{<fetched\_page>} block \\
\bottomrule
\end{tabularx}
\end{center}

The canonical trace files contain 100 Brave rows, 101 Tavily rows, and 100 Firecrawl rows. The extra Tavily row is a recovered transient failure; provider summaries select the latest valid row for each of the 100 query IDs, yielding the \NTraces{} analyzed provider-query trajectories.

\paragraph{Example trajectory.}
One qualitative failure case is \path{sealhard_a2c20fcc3ff1}: ``How many times did Blake Shelton win as a coach on The Voice (U.S.) between 2015 and his departure in 2023?'' The gold answer is 5, while all three provider runs answered 9. In the Brave trace \path{trace_5b6eea4e74d34191a4d275d656f82a18}, the agent used 2 iterations, made 1 search call, made 0 fetch calls, and answered from the pre-fetch surface. The provider-comparison row marks this as wrong despite answer text being visible somewhere in retrieved text. The point of the example is not that a single snippet is decisive, but that the pre-fetch surface mixed a period-specific answer with a tempting lifetime-win distractor.

\subsection{Trace Viewer and Qualitative Audit}
\label{app:trace-viewer}
The rendering script, \path{scripts/render_trace.py}, and local dashboard, \path{scripts/trace_dashboard.py}, support inspection of researcher-generated traces. We used rendered trajectory views to spot-check case studies, verify multi-call turns, and inspect failures where answer text was present but unused. These views are not an independent metric source and are excluded from the release when they reproduce provider content.

\section{Oracle and Metrics}
\label{app:metrics}
\label{app:judge-schema}
The Kimi-K2.6 judge takes one retrieved URL record as input. The exact judge prompt template is \path{src/searchapi_eval/evaluation/prompts/document_support_judge.liquid}; the execution wrapper is \path{scripts/run_llm_judge.py}. The judge receives the question, gold answer, model final answer, and one retrieved document rendered in the same XML-like representation that the answer model saw. Snippet-only records contain title, URL, domain, snippets, extra snippets, and metadata. Page-visible records additionally contain the fetched page block that was returned by \fetchPage{}.

Because the oracle is central to the paper's measurements, the table below states the guardrails used to keep each judge row constrained, parseable, and auditable.

\begin{center}
\footnotesize
\begin{tabularx}{\linewidth}{@{}L{0.34\linewidth}Z@{}}
\toprule
\textbf{Judge guardrail} & \textbf{Implementation} \\
\midrule
No external knowledge & Prompt restricts the judge to supplied document content. \\
Strict output & JSON-only schema with support, contradiction, garbage, spans, and confidence. \\
Valid-row filter & Requires schema version, parsed judgment, provider/query/retrieval/URL IDs, and no execution or parse error. \\
Garbage handling & Deterministic precheck flags failed, unsupported, empty, or media fetches before LLM judging. \\
Judgment reuse & Exact imports and question-URL reuse; see the deduplicated analysis below. \\
Error policy & Invalid judge rows are excluded, not converted into negative evidence. \\
\bottomrule
\end{tabularx}
\end{center}

The output schema fields used in the paper record gold support, snippet support, extracted-page support, model-answer support, contradiction, garbage status, evidence spans, and confidence. Across the main artifacts, \NJudgeValid{} of \NJudgeTotal{} rows are valid, split into \NJudgeSnippetValid{} snippet-only and \NJudgePageValid{} page-visible rows.

\subsection{Deduplicated Contradiction Analysis}
\label{app:cache-sensitivity}
The pooled analysis retains historical judgments, including question-URL reuse that did not require identical snippet text.

To assess repeated-URL weighting, we retain the first valid, independently judged snippet-only observation per provider, question, and normalized URL in canonical retrieval/result order. Exact cache imports qualify only when their full messages match the original request. Selection is independent of label value; all 5,634 retained observations have verified original-call provenance, and every eligible question-URL group is represented.

\begin{center}
\footnotesize
\setlength{\tabcolsep}{3pt}
\begin{tabularx}{\linewidth}{@{}L{0.22\linewidth}R{0.16\linewidth}R{0.16\linewidth}Y@{}}
\toprule
\textbf{Provider} & \textbf{Rows} & \textbf{c/g} & \textbf{Ratio [95\% CI]} \\
\midrule
Brave & 1,723 & 63/79 & 0.80 [0.44, 1.49] \\
Tavily & 2,059 & 54/28 & 1.93 [0.83, 5.00] \\
Firecrawl & 1,852 & 65/26 & 2.50 [1.07, 7.44] \\
\bottomrule
\end{tabularx}
\captionof{table}{Deduplicated contradiction-to-gold sensitivity. Rows are retained snippet-only observations; c/g gives contradiction/gold counts.}
\label{tab:app-deduplicated-ratios}
\end{center}

Intervals use 10,000 paired bootstrap samples of the 100 question IDs (seed 20260628), recomputing the ratio from retained observations in each sample. No replicate has a zero gold denominator. The paired differences are Brave--Tavily $-1.13$ [$-3.90$, $-0.17$], Brave--Firecrawl $-1.70$ [$-6.46$, $-0.44$], and Tavily--Firecrawl $-0.57$ [$-4.06$, $0.85$].

The derivation is implemented in \path{scripts/camera_ready_gate2_audit.py}.

\subsection{Metric Definitions}
For a provider $p$ and query $q$, let $U_{p,q}$ denote the valid judged URLs visible in the trajectory. Unfetched URLs have snippet-only judgments; fetched URLs have page-visible judgments that include both the pre-fetch snippet fields and the fetched page block. Let $G_{\mathrm{pre}}(u)$ be true when the pre-fetch surface supports the gold answer: either \goldSnip{} is true, or a snippet-only judgment marks \containsGold{} true. Let $G_{\mathrm{page}}(u)$ be true when a page-visible judgment marks \goldPage{} true. Let $F(u)$ be true when the trace records a \fetchPage{} attempt for $u$, including failed fetches and URLs without a valid judgment. Support predicates range over valid judged URLs; fetch existence ranges over all trace actions.

For each provider-query pair, pre-fetch support is $\exists u:G_{\mathrm{pre}}(u)$. Post-fetch discovered support is $\neg\exists u:G_{\mathrm{pre}}(u)$ and $\exists u:G_{\mathrm{page}}(u)$. Trajectory-visible support is their union.

The decision-cell table turns these definitions into the four labels used in the main paper. It is descriptive: it classifies what support was visible and what the agent fetched, rather than assigning causal blame.

\begin{center}
\footnotesize
\begin{tabularx}{\linewidth}{@{}L{0.15\linewidth}Z@{}}
\toprule
\textbf{Cell} & \textbf{Condition and interpretation} \\
\midrule
\smart{} & $\exists u:G_{\mathrm{pre}}(u)$ and $\exists u:G_{\mathrm{pre}}(u)\wedge F(u)$: the surface exposed support and the agent attempted to fetch it. \\
\missed{} & $\exists u:G_{\mathrm{pre}}(u)$ and $\neg\exists u:G_{\mathrm{pre}}(u)\wedge F(u)$: pre-fetch support was visible but not fetched. \\
\blind{} & $\neg\exists u:G_{\mathrm{pre}}(u)$ and $\exists u:F(u)$: the agent spent fetch budget without pre-fetch support. \\
\noop{} & $\neg\exists u:G_{\mathrm{pre}}(u)$ and $\neg\exists u:F(u)$: no pre-fetch support and no fetch. \\
\bottomrule
\end{tabularx}
\end{center}

The partition joins normalized URLs after resolving each fetch to its source retrieval and document ID. Fetch actions come directly from the canonical trace, independently of judge validity; invalid judgments remain excluded from support measurements. It is descriptive: a \missed{} query can still be answered correctly from snippets, and a \blind{} query can still succeed if a fetched page reveals evidence not present in snippets. The pooled surface contradiction-to-gold ratio counts snippet-only URL observations, including repeated appearances:
\[
 r_{c:g}=\frac{|\{u:\contraGold(u)\}|}{|\{u:\containsGold(u)\}|}.
\]
It is independent of the model's final answer, but not independent of the gold answer.

For answer-level scoring, exact match and token F1 use the same deterministic normalizer: lowercasing, removing articles and punctuation, and mapping simple number words to digits. Token F1 is the token-overlap F1 between the normalized final answer and normalized gold answer, macro-averaged over the \NQueries{} queries. The semantic \correctans{} metric is separate and comes from the audited \semanticMatch{} labels in \path{results/released/answer_audit.csv}.

\subsection{Human Validation and Semantic Corrections}
\label{app:examples}
The validation queue samples row/label pairs using \path{scripts/build_judge_validation_queue.py}; \path{scripts/judge_validation_app.py} records one human decision per case. One author conducted the review with provider identity hidden. The completed CSV contains 180 cases: five per provider $\times$ surface (snippet-only/page-visible) $\times$ label $\times$ Kimi positive/negative value. The audited labels are \containsGold{}, \contraGold{}, and \texttt{is\_garbage}; \goldSnip{} and \goldPage{} were not directly audited.

The reviewer recorded true, false, or unclear. Six unclear decisions remain in the sample and are excluded only from the agreement denominator. The 174 clear decisions include 164 agreements and 10 disagreements. The table describes this balanced audit sample, not a population-wide error rate or agreement between independent annotators.

\begin{center}
\footnotesize
\begin{tabularx}{\linewidth}{@{}L{0.34\linewidth}R{0.13\linewidth}R{0.13\linewidth}R{0.22\linewidth}@{}}
\toprule
\textbf{Slice} & \textbf{Clear} & \textbf{Agree} & \textbf{Agreement} \\
\midrule
Gold support & 58 & 55 & 94.8\% \\
Contradiction & 56 & 53 & 94.6\% \\
Garbage & 60 & 56 & 93.3\% \\
Overall & 174 & 164 & 94.3\% \\
\bottomrule
\end{tabularx}
\end{center}

For the separate final-answer audit, exact matches were automatically accepted and one author reviewed non-exact answers with provider identity hidden. The audit covers all 300 outcomes: 65 exact matches and 11 additional semantic matches, yielding 25/25/26 correct answers for Brave/Tavily/Firecrawl. Accepted corrections include spacing (\emph{Astra Zeneca}/\emph{AstraZeneca}), entity naming (\emph{UnionPay}/\emph{China UnionPay}), and equivalent numeric expressions (\emph{3}/\emph{3 players}). Recorded notes retain ambiguous specificity (\emph{Punjabi}/\emph{Western Punjabi}), rounded values, hedged ranges, and over-inclusive answers as non-matches. The released CSV preserves the individual labels; no answer labels were changed for the camera-ready version.

\section{Results}
\label{app:additional}
The result appendices report uncertainty, decision-cell diagnostics, token and fetch diagnostics, and artifact provenance. These tables support the main claims without changing the headline conclusions.

\subsection{Query-Level Diagnostics}
\label{app:query-diagnostics}
First-search metrics use the first recorded search call in each canonical provider--question trajectory, before later query reformulation. Support is defined by $G_{\mathrm{pre}}$ in Appendix~\ref{app:metrics}; page-only support does not contribute. Support@1 is the proportion of all 100 questions with observed support at rank 1. Mean reciprocal rank averages $1/r$ for the first observed supporting result, with zero when no support is observed. Support at any rank also uses all 100 questions.

Invalid judgments remain unknown. The table reports their coverage and deterministic lower/upper bounds obtained by completing all missing labels as negative/positive, while holding valid labels fixed. These bounds are separate from the bootstrap intervals. Firecrawl has one empty first-search response, which contributes zero rather than a missing judgment.

\begin{center}
\footnotesize
\setlength{\tabcolsep}{3pt}
\begin{tabularx}{\linewidth}{@{}L{0.42\linewidth}YYY@{}}
\toprule
\textbf{Diagnostic} & \textbf{Brave} & \textbf{Tavily} & \textbf{Firecrawl} \\
\midrule
Valid/returned first-search labels & 995/1000 & 958/969 & 978/990 \\
Unknown rank-1 labels & 0 & 5 & 3 \\
Median first-support rank (supported questions) & 3 ($n=24$) & 1 ($n=12$) & 3 ($n=9$) \\
\midrule
Support@1 bounds (\%) & 5--5 & 9--14 & 2--5 \\
Mean reciprocal rank bounds & 0.103--0.107 & 0.103--0.157 & 0.041--0.079 \\
Support-at-any-rank bounds (\%) & 24--26 & 12--19 & 9--17 \\
Contradiction-exposure bounds (\%) & 27--32 & 21--26 & 28--36 \\
\bottomrule
\end{tabularx}
\captionof{table}{First-search coverage, conditional first-support rank, and missing-label bounds. Contradiction exposure uses snippet-only observations over the full trajectory. Bounds concern missing judgments; they are not confidence intervals.}
\label{tab:app-query-coverage}
\end{center}

Contradiction diagnostics use valid snippet-only observations over the full trajectory, counting repeated appearances. Page-visible contradiction labels may reflect page content, so those observations are excluded from this snippet-only analysis. Every question has at least one valid snippet-only observation. Brave/Tavily/Firecrawl have 89/58/70 contradictory observations out of 2,095/2,339/2,085 valid observations, with 97/31/27 gold-supporting observations and 6/12/14 invalid observations. The independent-observation selection in Appendix~\ref{app:cache-sensitivity} leaves question-exposure counts at 27/21/28; the corresponding contradictory-observation percentages are 3.66/2.62/3.51.

Intervals use 10,000 matched-question percentile-bootstrap samples, seed 20260628, recomputing ratios from observation counts in each sample. For Brave--Tavily, Brave--Firecrawl, and Tavily--Firecrawl, respectively, Support@1 differences are $-4$ [$-11$, 3], 3 [0, 7], and 7 [2, 13] percentage points; MRR differences are $-0.0004$ [$-0.067$, 0.064], 0.062 [0.023, 0.107], and 0.063 [0.008, 0.124]; contradiction-exposure differences are 6 [$-3$, 15], $-1$ [$-10$, 8], and $-7$ [$-15$, 1] percentage points. These intervals describe the observed metrics; the missing-label bounds above address incomplete judgments separately.

The first-search query strings are identical across all three providers for 74 questions. Within this subset, observed Support@1 is 4/74, 8/74, and 2/74, support at any rank is 20/74, 10/74, and 7/74, and MRR is 0.117, 0.119, and 0.047 (Brave/Tavily/Firecrawl). This descriptive check retains the distinction between support availability and position. Per-question results, the full first-support-rank distribution, and paired intervals from \path{scripts/camera_ready_gate3_diagnostics.py} are released under \path{results/released/}.

\subsection{Paired Bootstrap Uncertainty}
\label{app:uncertainty}
We resample question IDs with replacement, preserving the matched Brave/Tavily/Firecrawl rows for each question. The tables use 10,000 percentile-bootstrap replicates with seed 20260628. Count metrics are reported per 100 sampled questions; pairwise differences are left minus right in the row's native units.

The first uncertainty table gives provider-level intervals for the main quantities. These intervals show the scale of each provider's measured regime before any pairwise comparison is made. The rank-1 row uses the same denominator as Table~\ref{tab:visible-support}: gold-supporting pre-fetch rows.

\begin{center}
\footnotesize
\setlength{\tabcolsep}{3pt}
\begin{tabularx}{\linewidth}{@{}L{0.34\linewidth}Z@{}}
\toprule
\textbf{Metric} & \textbf{Provider intervals} \\
\midrule
\correctans{} /100 & B 25 [17, 34]; T 25 [17, 34]; F 26 [18, 35] \\
Pre-fetch support /100 & B 30 [21, 39]; T 16 [9, 23]; F 16 [9, 23] \\
Rank-1 among pre-fetch support rows & B 12.9 [7.1, 19.4]; T 50.0 [31.2, 78.9]; F 13.3 [4.2, 26.3] \\
$r_{c:g}$ & B 0.92 [0.49, 1.73]; T 1.87 [0.83, 4.47]; F 2.59 [1.11, 8.00] \\
Fetched queries /100 & B 65 [56, 74]; T 76 [67, 84]; F 81 [73, 89] \\
Avg. fetch calls & B 1.02 [0.81, 1.25]; T 1.30 [1.06, 1.55]; F 1.28 [1.08, 1.50] \\
Tokens/query & B 59,627 [47,706, 72,908]; T 54,156 [44,304, 64,991]; F 57,979 [44,635, 72,878] \\
\bottomrule
\end{tabularx}
\captionof{table}{Provider bootstrap 95\% intervals. B=Brave, T=Tavily, F=Firecrawl.}
\label{tab:app-bootstrap-provider}
\end{center}

The second uncertainty table reports matched pairwise differences. This comparison clarifies the main claim boundary: intervals for correctness differences include zero and do not establish equivalence, while evidence-economy contrasts are more stable.

\begin{center}
\footnotesize
\setlength{\tabcolsep}{3pt}
\begin{tabularx}{\linewidth}{@{}L{0.34\linewidth}Z@{}}
\toprule
\textbf{Metric} & \textbf{Pairwise differences} \\
\midrule
\correctans{} /100 & B--T 0 [$-9$, 9]; B--F $-1$ [$-10$, 8]; T--F $-1$ [$-10$, 8] \\
Pre-fetch support /100 & B--T 14 [4, 24]; B--F 14 [6, 22]; T--F 0 [$-8$, 8] \\
Rank-1 among pre-fetch support rows & B--T $-37.1$ [$-67.3$, $-16.8$]; B--F $-0.5$ [$-13.3$, 10.1]; T--F 36.7 [18.8, 66.4] \\
$r_{c:g}$ & B--T $-0.95$ [$-3.22$, 0.00]; B--F $-1.68$ [$-6.78$, $-0.36$]; T--F $-0.72$ [$-4.86$, 0.53] \\
Fetched queries /100 & B--T $-11$ [$-19$, $-3$]; B--F $-16$ [$-25$, $-7$]; T--F $-5$ [$-12$, 1] \\
Avg. fetch calls & B--T $-0.28$ [$-0.55$, $-0.02$]; B--F $-0.26$ [$-0.51$, $-0.02$]; T--F 0.02 [$-0.23$, 0.27] \\
Tokens/query & B--T 5,471 [$-4{,}591$, 16,176]; B--F 1,648 [$-11{,}834$, 13,959]; T--F $-3{,}823$ [$-18{,}019$, 9,064] \\
\bottomrule
\end{tabularx}
\captionof{table}{Paired bootstrap differences; left minus right.}
\label{tab:app-bootstrap-diff}
\end{center}

\subsection{Additional Diagnostics and Artifact Manifest}
The remaining diagnostics are lower-level checks behind the headline tables. They are separated from the main results because they support interpretation but are not themselves the paper's primary claims.

The retrieval-proxy table reports deterministic string and URL checks from traces. These proxies help show why answer visibility is not the same as correct evidence use.

\begin{center}
\footnotesize
\setlength{\tabcolsep}{2pt}
\begin{tabularx}{\linewidth}{@{}L{0.45\linewidth}Y@{}}
\toprule
\textbf{Metric} & \textbf{B/T/F} \\
\midrule
Gold URL exact hit & \BraveGoldURLExact{} / \TavilyGoldURLExact{} / \FirecrawlGoldURLExact{} \\
Gold domain hit & \BraveGoldDomain{} / \TavilyGoldDomain{} / \FirecrawlGoldDomain{} \\
Gold source-family hit & \BraveGoldFamily{} / \TavilyGoldFamily{} / \FirecrawlGoldFamily{} \\
Answer in snippet surface & \BraveSnippetSurfaceHit{} / \TavilySnippetSurfaceHit{} / \FirecrawlSnippetSurfaceHit{} \\
Answer in fetched page & \BraveAnswerPage{} / \TavilyAnswerPage{} / \FirecrawlAnswerPage{} \\
Answer in any retrieved text & \BraveAnswerAvailable{} / \TavilyAnswerAvailable{} / \FirecrawlAnswerAvailable{} \\
Wrong despite answer text & \BraveWrongWithAnswer{} / \TavilyWrongWithAnswer{} / \FirecrawlWrongWithAnswer{} \\
\bottomrule
\end{tabularx}
\captionof{table}{Trace-derived retrieval proxies. B/T/F = Brave/Tavily/Firecrawl; these deterministic URL and answer-string checks are not oracle support labels.}
\label{tab:app-retrieval-diagnostics}
\end{center}

The Wilson-interval table expands the decision partition by provider and cell. It is included because some cells are small, so raw cell accuracies should be read with uncertainty.

\begin{center}
\begin{minipage}{\linewidth}
\footnotesize
\begin{tabularx}{\linewidth}{@{}L{0.34\linewidth}Z@{}}
\toprule
\textbf{Provider/cell} & \textbf{Correct rate and Wilson 95\% CI} \\
\midrule
Brave \smart{} & \BraveSmartCorrect{}/\BraveSmartN{}: \BraveSmartRate{} \BraveSmartCI{} \\
Brave \missed{} & \BraveMissedCorrect{}/\BraveMissedN{}: \BraveMissedRate{} \BraveMissedCI{} \\
Brave \blind{} & \BraveBlindCorrect{}/\BraveBlindN{}: \BraveBlindRate{} \BraveBlindCI{} \\
Brave \noop{} & \BraveNoopCorrect{}/\BraveNoopN{}: \BraveNoopRate{} \BraveNoopCI{} \\
Tavily \smart{} & \TavilySmartCorrect{}/\TavilySmartN{}: \TavilySmartRate{} \TavilySmartCI{} \\
Tavily \missed{} & \TavilyMissedCorrect{}/\TavilyMissedN{}: \TavilyMissedRate{} \TavilyMissedCI{} \\
Tavily \blind{} & \TavilyBlindCorrect{}/\TavilyBlindN{}: \TavilyBlindRate{} \TavilyBlindCI{} \\
Tavily \noop{} & \TavilyNoopCorrect{}/\TavilyNoopN{}: \TavilyNoopRate{} \TavilyNoopCI{} \\
Firecrawl \smart{} & \FirecrawlSmartCorrect{}/\FirecrawlSmartN{}: \FirecrawlSmartRate{} \FirecrawlSmartCI{} \\
Firecrawl \missed{} & \FirecrawlMissedCorrect{}/\FirecrawlMissedN{}: \FirecrawlMissedRate{} \FirecrawlMissedCI{} \\
Firecrawl \blind{} & \FirecrawlBlindCorrect{}/\FirecrawlBlindN{}: \FirecrawlBlindRate{} \FirecrawlBlindCI{} \\
Firecrawl \noop{} & \FirecrawlNoopCorrect{}/\FirecrawlNoopN{}: \FirecrawlNoopRate{} \FirecrawlNoopCI{} \\
\bottomrule
\end{tabularx}
\captionof{table}{Decision-cell Wilson 95\% intervals.}
\label{tab:app-wilson}
\end{minipage}
\end{center}

The overlap table gives a compact view of which providers answer the same questions correctly. It supports the complementarity claim by separating shared successes from provider-specific successes.

\begin{center}
\footnotesize
\setlength{\tabcolsep}{2pt}
\begin{tabularx}{\linewidth}{@{}L{0.45\linewidth}Y@{}}
\toprule
\textbf{Pair} & \textbf{both / left / right / wrong} \\
\midrule
Brave vs. Firecrawl & 14 / 11 / 12 / 63 \\
Brave vs. Tavily & 14 / 11 / 11 / 64 \\
Firecrawl vs. Tavily & 14 / 12 / 11 / 63 \\
\bottomrule
\end{tabularx}
\captionof{table}{Semantic-correct pairwise overlap: both / left / right / wrong.}
\label{tab:app-pairwise}
\end{center}

The final diagnostic table reports legacy token F1 and token and fetch volume. These counts ground the paper's cost and retrieval-budget discussion in the actual traced runs.

\begin{center}
\footnotesize
\setlength{\tabcolsep}{2pt}
\begin{tabularx}{\linewidth}{@{}L{0.45\linewidth}Y@{}}
\toprule
\textbf{Metric} & \textbf{B/T/F} \\
\midrule
Token F1 & \BraveFone{} / \TavilyFone{} / \FirecrawlFone{} \\
Total tokens & \BraveTokensM{}M / \TavilyTokensM{}M / \FirecrawlTokensM{}M \\
Median/query & \BraveMedianTokens{} / \TavilyMedianTokens{} / \FirecrawlMedianTokens{} \\
Max/query & \BraveMaxTokens{} / \TavilyMaxTokens{} / \FirecrawlMaxTokens{} \\
$>$100k queries & \BraveOverHundredK{} / \TavilyOverHundredK{} / \FirecrawlOverHundredK{} \\
Fetch success/fail & \BraveFetchSuccess{}/\BraveFetchFailed{} / \TavilyFetchSuccess{}/\TavilyFetchFailed{} / \FirecrawlFetchSuccess{}/\FirecrawlFetchFailed{} \\
\bottomrule
\end{tabularx}
\captionof{table}{Legacy token F1, token usage, and fetch diagnostics. B/T/F = Brave/Tavily/Firecrawl.}
\label{tab:app-token}
\end{center}

\begin{description}[leftmargin=0pt,itemsep=0.12em,topsep=0.1em]
\item[Analysis inputs.] The reported analyses used the semantic-answer audit, three provider trace files, three Kimi per-URL judgment files, and provider-comparison summaries.
\item[Released artifacts.] The release includes code, prompt templates, configurations, sampling procedures, schemas, evaluation and visualization scripts, paper sources, and derived results without copied provider content. Complete traces and content-bearing judge, audit, and trajectory-view records are excluded.
\end{description}
